% Full answers PDF: MCQ answers + worked explanations per option, then official MS clips.
% Compiler: XeLaTeX.
\documentclass[12pt,a4paper]{article}
\usepackage[margin=18mm,headheight=14pt]{geometry}
\usepackage{fontspec}
\usepackage[version=4]{mhchem}
\usepackage{graphicx}
\usepackage{xcolor}
\usepackage{booktabs}
\usepackage{fancyhdr}
\usepackage{enumitem}
\usepackage{needspace}

\setmainfont{Helvetica Neue}

\pagestyle{fancy}
\fancyhf{}
\fancyhead[L]{\small 3.1 Electronegativity --- answers}
\fancyhead[R]{\small CIE 9701}
\fancyfoot[C]{\thepage}
\renewcommand{\headrulewidth}{0.3pt}

\setlength{\parindent}{0pt}
\setlength{\parskip}{0.35em}

% ---- helper: mark-scheme image block ----
\newcommand{\msblock}[2]{%
  \par\vspace{0.8em}%
  {\normalsize\bfseries #1}\par\vspace{0.25em}%
  \noindent\includegraphics[width=\textwidth,keepaspectratio]{#2}\par
}

% ---- helper: one MCQ worked explanation ----
% \mcq{<question>}{<correct letter>}{<source>}{<body>}
\newcommand{\mcq}[4]{%
  \par\needspace{5\baselineskip}\vspace{0.7em}%
  {\normalsize\bfseries #1.\quad #2}%
  \hfill{\footnotesize\color{gray}#3}\par\vspace{0.15em}%
  {\small #4}\par
}

% ---- helper: a "distractor" line ----
\newcommand{\why}[2]{\par\needspace{2.2\baselineskip}\vspace{0.1em}{\small\textbf{#1} #2}\par}

\begin{document}

\begin{center}
{\large\bfseries 3.1 Electronegativity homework --- answers}\\[0.2em]
{\small CIE 9701 AS Chemistry \quad·\quad 11 multiple-choice + 5 written marks}
\end{center}

{\small\bfseries Do not issue this file to students.}
\hrule
\vspace{0.4em}

% =====================================================================
\section*{Section A --- Multiple choice}
% =====================================================================

{\normalsize\bfseries Quick answers}

\vspace{0.3em}
{\small
\begin{tabular}{@{}ll@{\hspace{2.6em}}ll@{\hspace{2.6em}}ll@{}}
\toprule
Q & Ans. & Q & Ans. & Q & Ans.\\
\midrule
A1 & \textbf{A} & A5 & \textbf{B} & A9 & \textbf{C}\\
A2 & \textbf{C} & A6 & \textbf{B} & A10 & \textbf{B}\\
A3 & \textbf{C} & A7 & \textbf{D} & A11 & \textbf{C}\\
A4 & \textbf{C} & A8 & \textbf{A} & & \\
\bottomrule
\end{tabular}}

\vspace{0.6em}
{\normalsize\bfseries Worked explanations}

\mcq{A1}{A}{2022 MJ Paper 12 Q19}{%
Across Period 3 from Na to Si the nuclear charge rises while the outer electrons
stay in the same shell, so the atomic radius \emph{decreases} and the
electronegativity \emph{increases}. The correct graph must therefore slope
downwards, with Si (highest electronegativity, smallest radius) at the top-left
and Na (lowest electronegativity, largest radius) at the bottom-right.}

\why{B:}{the elements are in the right order on the radius axis, but Al is placed
above Si --- Si has the highest electronegativity of the four, so Si must be the
highest point.}

\why{C:}{shows electronegativity \emph{increasing} with atomic radius (Na at the
top, Si at the bottom) --- the trend is the wrong way round.}

\why{D:}{swaps Al and Si on the radius axis (Al appears to the left of Si), but
the correct radius order is Si $<$ Al $<$ Mg $<$ Na.}

\mcq{A2}{C}{2023 FM Paper 12 Q22}{%
Use the two radii to identify the elements.
V has a large atomic radius (0.186 nm) but a much smaller ionic radius
(0.095 nm) --- it is a metal that has \emph{lost} its outer electron. This matches
Na (Na 0.186 nm, \ce{Na+} 0.095 nm).
Z has a small atomic radius (0.099 nm) but a \emph{larger} ionic radius
(0.181 nm) --- it is a non-metal that has \emph{gained} an electron. This matches
Cl (Cl 0.099 nm, \ce{Cl-} 0.181 nm). So V = Na and Z = Cl.}

\why{A:}{\ce{Na+} has two shells (2,8) but \ce{Cl-} has three (2,8,8), so they do
\emph{not} have the same number of full shells.}

\why{B:}{Z is chlorine, which forms a \emph{negative} ion, \ce{Cl-}.}

\why{C (correct):}{Cl is far more electronegative than Na (3.2 vs 0.9).}

\why{D:}{Na has 1 outer electron and Cl has 7, so V has \emph{fewer} outer
electrons than Z.}

\mcq{A3}{C}{2022 FM Paper 12 Q18}{%
Electronegativity \emph{increases} across Period 3 (nuclear charge increases,
shielding stays roughly constant, radius falls). The oxides change from
\emph{basic} (\ce{Na2O}, MgO) through \emph{amphoteric} (\ce{Al2O3}) to
\emph{acidic} (\ce{SiO2}, \ce{P4O10}, \ce{SO3}, \ce{Cl2O7}). Only row C has both
trends correct.}

\why{A / B:}{both state that electronegativity \emph{decreases} across the
period --- the wrong direction.}

\why{D:}{has the correct electronegativity trend but the acid--base order
reversed.}

\mcq{A4}{C}{2023 MJ Paper 13 Q17}{%
Astatine is below iodine in Group 17. Going down the group the atomic radius and
the shielding increase, so the nucleus attracts the bonding electrons less and
the electronegativity \emph{decreases}. The larger, more electron-rich molecule
also has stronger van der Waals' forces between molecules, so astatine is
\emph{less volatile} (higher boiling point) than iodine.}

\why{A / B:}{both say astatine is \emph{more} electronegative --- wrong for a
heavier halogen.}

\why{D:}{correct lower electronegativity, but says astatine is \emph{more}
volatile; volatility decreases down the group.}

\mcq{A5}{B}{2025 MJ Paper 11 Q20}{%
Electronegativity increases across Na $\rightarrow$ Mg $\rightarrow$ Al
$\rightarrow$ Si. First ionisation energy also increases \emph{overall}, but not
in atomic-number order: Mg (3s\textsuperscript{2}, a full subshell) has a
\emph{higher} first ionisation energy than Al (3p\textsuperscript{1}, easier to
remove). So the graph must show the points in the order
Na $<$ Al $<$ Mg $<$ Si along the ionisation-energy axis, while the
electronegativity order is Na $<$ Mg $<$ Al $<$ Si --- i.e.\ Mg lies above Al,
and Si is both furthest right and highest.}

\why{A:}{places Mg before Al on the ionisation-energy axis, giving Al a higher
first ionisation energy than Mg --- this ignores the 3s\textsuperscript{2} /
3p\textsuperscript{1} difference.}

\why{C / D:}{show electronegativity \emph{decreasing} as ionisation energy
increases; in fact the two trends rise together across the period.}

\mcq{A6}{B}{2025 ON Paper 14 Q17}{%
Electronegativity increases Mg $<$ Al $<$ Si $<$ P, so the points must appear in
that order from left to right. The melting point depends on structure: Mg and Al
are metals with similar melting points ($\approx$ 650--660\,\textdegree C); Si is
a giant covalent (macromolecular) solid with a much higher melting point
(1414\,\textdegree C); P exists as simple \ce{P4} molecules held by weak van der
Waals' forces, so it melts at only 44\,\textdegree C. Graph B is the only one
showing Si clearly the highest, Mg and Al close together and much lower, and P
the lowest.}

\why{A:}{puts Al above Si --- the melting point of Si is far higher than that of
Al.}

\why{C / D:}{place the elements in the wrong order on the electronegativity axis
(P to the left of Si/Mg), reversing the Period 3 trend.}

\mcq{A7}{D}{2022 MJ Paper 13 Q5}{%
The difference in electronegativity is smallest for the pair of elements that are
closest in electronegativity. Using Pauling values
(K 0.8, Li 1.0, Br 3.0, F 4.0):
$\Delta E_\mathrm{N}(\ce{KF}) \approx 3.2$,
$\Delta E_\mathrm{N}(\ce{LiF}) \approx 3.0$,
$\Delta E_\mathrm{N}(\ce{KBr}) \approx 2.1$,
$\Delta E_\mathrm{N}(\ce{LiBr}) \approx 2.0$.
Li and Br are the closest pair, so \ce{LiBr} has the smallest difference.}

\why{A / C:}{fluorine is the most electronegative element, so both \ce{KF} and
\ce{LiF} give large differences.}

\why{B:}{\ce{KBr} has the same anion as D, but K is less electronegative than Li,
so its difference is larger.}

\mcq{A8}{A}{2023 MJ Paper 13 Q4}{%
L has the \emph{highest} Pauling value of the seven elements Na--Cl, so L is
chlorine. M has the \emph{lowest}, so M is sodium.
(1) Chlorine exists as \ce{Cl2}, which contains a covalent Cl--Cl bond ---
\textbf{true}.
(2) Cl ($Z = 17$) has a higher atomic number than Na ($Z = 11$) ---
\textbf{true}.
(3) Sodium and chlorine form \ce{NaCl}, a giant ionic lattice --- \textbf{true}.
All three statements are correct, so the answer is A.}

\why{B / C / D:}{each option rejects at least one statement that is in fact
correct --- the trap is assuming that an element cannot contain covalent bonds
(it can, e.g.\ \ce{Cl2}), or that Na and Cl would bond covalently.}

\mcq{A9}{C}{2022 FM Paper 12 Q7}{%
Identify the elements from their electronegativity: X (1.0) = Na, Y (2.1) = H,
Z (4.0) = F.
\ce{Z2} = \ce{F2}: simple molecular, only weak van der Waals' forces ---
\emph{lowest} melting point.
\ce{YZ} = HF: simple molecular with hydrogen bonding --- higher than \ce{F2}, but
still low.
\ce{XZ} = NaF: giant ionic lattice with strong electrostatic forces ---
\emph{highest} melting point.
So increasing melting point is \ce{Z2} $<$ \ce{YZ} $<$ \ce{XZ}, which is row C.}

\why{A / B:}{put the ionic compound \ce{XZ} at the lowest end --- ionic lattices
have by far the strongest forces here.}

\why{D:}{puts \ce{YZ} (HF) highest; hydrogen bonding between small molecules is
much weaker than the electrostatic forces in an ionic lattice.}

\mcq{A10}{B}{2021 MJ Paper 13 Q6}{%
Ionic bonding requires the \emph{transfer} of electrons from a metal atom to a
non-metal atom, i.e.\ two different elements. A single element contains only one
type of atom and cannot form oppositely charged ions, so ionic bonding is
\emph{never} found in elements.}

\why{A:}{covalent --- \ce{H2}, \ce{Cl2}, diamond.}

\why{C:}{metallic --- Na, Fe.}

\why{D:}{van der Waals' forces --- between \ce{I2} molecules or noble-gas atoms.}

\mcq{A11}{C}{2023 MJ Paper 12 Q19}{%
(1) \textbf{True} --- the ionically bonded oxides \ce{Na2O} and MgO are basic and
neutralise dilute hydrochloric acid.
(2) \textbf{False} --- this is the trap. \ce{NaCl} does give a neutral solution,
but \ce{MgCl2} gives an \emph{acidic} solution (the small, highly charged
\ce{Mg^2+} ion polarises water), so not all metal chlorides are neutral.
(3) \textbf{True} --- the two most electronegative elements of the set are Si and
P, and \ce{SiCl4} and \ce{PCl3}/\ce{PCl5} are covalent.
Correct statements: 1 and 3 $\rightarrow$ C.}

\newpage

% =====================================================================
\section*{Section B --- Mark scheme clips}
% =====================================================================

\msblock{B1(a)(i)  [1]}{cie-9701-2023-fm-p22-q1a-i-ms.png}

\msblock{B1(a)(ii)  [2]}{cie-9701-2023-fm-p22-q1a-ii-ms.png}

\msblock{B2(b)(ii)  [2]}{cie-9701-2023-mj-p22-q1b-ii-ms.png}

\end{document}
